\section{Deployment profiles and configuration}
\label{supp:deployment}

This supplementary section enumerates the v2 deployment profiles, the
configuration keys that change between profiles, and the runtime
constraints that apply to each profile.  The profiles are mutually
exclusive; a deployment that mixes two profiles is misconfigured.

\subsection{Single-tenant, single-process (S1P)}

\paragraph{Use case.}  Local research, bundled benchmarks, a single
researcher running one world.

\paragraph{Configuration.}
\texttt{accounts.enabled = false} (the default).  The
\texttt{GAWORLD\_DASHBOARD\_TOKEN} environment variable is the only
credential; a Bearer header in an HTTP request authenticates as admin.

\paragraph{Runtime constraints.}  At most one running simulation.
All paths under \texttt{output/} are used directly.  No per-tenant
quotas apply.

\subsection{Multi-tenant, single-process (M1P)}

\paragraph{Use case.}  A classroom, a workshop, a small research group
sharing a single server.

\paragraph{Configuration.}
\texttt{accounts.enabled = true}; the account store is initialized
with \texttt{python -m gaworld.accounts init --admin <nickname>}.
Invites are issued in batches with \texttt{python -m gaworld.accounts
invite --count N --label <class>}.  Quotas are configured by the
admin through \texttt{POST /api/worlds/settings}.

\paragraph{Runtime constraints.}
A \emph{default} world always exists for back-compat; non-default
worlds are stored under \texttt{output/worlds/<id>/}.  A world has
at most one running simulation; the deployment as a whole has at
most \texttt{max\_concurrent\_runs} (default 4).  Each user has at
most \texttt{max\_runs\_per\_user} (default 1) running simulations;
admins are exempt.

\subsection{Multi-tenant, multi-node (M1N)}

\paragraph{Use case.}  A school lab, a cross-campus collaboration, a
multi-region test, a teacher whose students run their own model keys
on their own laptops.

\paragraph{Configuration.}
Same as M1P, plus at least one node registered with
\texttt{POST /api/worlds/<id>/nodes} and joined with
\texttt{python -m gaworld.cluster join <hub-url> --token <token>}.

\paragraph{Runtime constraints.}
Same as M1P, plus a per-node intervention poll interval (default 1 s)
and a per-node record-upload interval (default 1 s).  The hub
publishes a logical clock that is the median of the per-node clocks;
the visualizer and the intervention display consult this clock but
do not require it for correctness.  A node that disconnects keeps
running locally and back-fills records on reconnection; a node
whose token is revoked receives 401 on its next request and stops
the simulation.

\subsection{Backwards compatibility}

The \texttt{AccountStore} is optional.  An installation that has not
run \texttt{python -m gaworld.accounts init} continues to behave
exactly as v1 did, with \texttt{GAWORLD\_DASHBOARD\_TOKEN} as the only
credential and a single \texttt{default} world.  Migration to
multi-tenancy is a one-way operation: the account store cannot be
removed without deleting every world.  The v2 path-resolution layer
honors the v1 module constants when \texttt{accounts.enabled = false},
so unit tests that monkey-patch module constants continue to work.

\subsection{Hardening checklist}

For a production deployment, the following additional hardening steps
are recommended.  They are not part of v2 by default and are not
required for local research.

\begin{itemize}[leftmargin=*]
  \item \textbf{TLS.}  Terminate TLS at a reverse proxy (nginx,
  Caddy, or a managed load balancer).  Set
  \texttt{HSTS}, \texttt{Secure} cookie, and a redirect from HTTP
  to HTTPS.
  \item \textbf{Authentication.}  Replace
  \texttt{GAWORLD\_DASHBOARD\_TOKEN} with a per-user token issued by
  the account store.  Enable rate limiting and brute-force
  protection on \texttt{/api/auth}.
  \item \textbf{Authorization.}  Audit the per-role permissions in
  \texttt{gaworld/accounts/policy.py} and confirm that the role
  \texttt{owner} cannot exceed \texttt{max\_concurrent\_runs} by
  hiding as \texttt{admin}.
  \item \textbf{Quota.}  Configure a per-tenant LLM-call quota
  appropriate to the underlying provider's billing model.  A dry-run
  mode logs the would-be denial without refusing the request.
  \item \textbf{Audit.}  Forward the \texttt{audit} table to a
  long-term store (a SQL database or a log-shipping service) and
  retain the records for the period required by the institution.
  \item \textbf{Redis / database.}  Replace the in-process
  \texttt{DistributedRelayBackend} with a Redis-backed relay if
  the deployment crosses a process boundary that the in-process
  relay cannot serve.  Replace the SQLite \texttt{accounts.sqlite}
  with a managed database if the deployment has more than one
  dashboard process.
  \item \textbf{Secret management.}  Load \texttt{OPENAI\_API\_KEY},
  \texttt{ANTHROPIC\_API\_KEY}, and other provider credentials from
  a secret manager (AWS Secrets Manager, GCP Secret Manager, or
  HashiCorp Vault) rather than from a dotenv file in the working
  directory.
\end{itemize}

These steps do not change the protocol semantics; they raise the
deployment to a production bar.  Without them, a public deployment
exposes the research surface to the open Internet and inherits the
risks of any such exposure.
